%=============================================================================!
% 2.F  SOLUTIONS TO THE EXERCISES                                             !
%=============================================================================!
\unnumbered{Chapter~\ref{ch:newton}: Forces and Newton's laws}
\label{sec:ne-solutions}

\solutionto{ex:ne-lift}The passenger moves at constant velocity, so by
the first law the net force on them is zero: the floor pushes up exactly
as hard as their weight pulls down. A frame of reference fixed to the
lift moves at constant velocity relative to the ground, which is very
nearly inertial, so by \cref{sec:ne-frames} it is very nearly inertial
too. While the
lift speeds up, it accelerates, and a frame of reference fixed to it is
no longer inertial. Seen from the ground, the passenger then accelerates
upwards with the lift, so the net force on them points up: the floor
pushes harder than their weight, which is why they feel heavier.

\solutionto{ex:ne-frames}By \cref{eq:ne-velocity-frames}, $v' = v - u$,
and adding $u$ to both sides, the velocity seen from the ground is $v =
v' + u$, the velocity seen from the bicycle plus the velocity of the
bicycle. Taking forwards as positive, $u =
\SI{6}{\metre\per\second}$. Thrown forwards, $v' = +4$ and $v = 6 + 4 =
\SI{10}{\metre\per\second}$. Thrown backwards, $v' = -4$ and $v = 6 - 4 =
\SI{2}{\metre\per\second}$: the ball still moves forwards relative to the
ground, only more slowly than the bicycle.

\solutionto{ex:ne-net}Taking east as $x$ and north as $y$, the net force
is $\vb{F} = (4, 0) + (0, 3) = (4, 3)$\,\si{\newton}. Its size is $(4^2 +
3^2)^{1/2} = \SI{5}{\newton}$, and its direction makes an angle $\theta$
north of east with $\tan\theta = 3/4$, so $\theta = \arctan(3/4) =
\ang{36.9}$. By the second law the acceleration points the same way and
has size $F/m = 5/3 = \SI{1.67}{\metre\per\second\squared}$.

\solutionto{ex:ne-mass}The same force gives accelerations in the inverse
ratio of the masses, $m_2/m_1 = a_1/a_2$, so
\begin{equation*}
   m_2 = m_1\,\frac{a_1}{a_2} = 2\times\frac{0.6}{1.5}
   = \SI{0.8}{\kilogram} .
\end{equation*}
The force follows from either trolley by the second law: $F = m_1 a_1 =
2\times0.6 = \SI{1.2}{\newton}$, and as a check $m_2 a_2 = 0.8\times1.5 =
\SI{1.2}{\newton}$.

\solutionto{ex:ne-horse}The free-body diagrams below draw the cart and
the horse apart, moving to the right. The weight and normal force on each
balance, the teal arrows are the third-law pair between them, and the
lengths of the horizontal arrows show which forces are larger.
\begin{center}
\begin{tikzpicture}[line cap=round,
   force/.style={-{Stealth[length=5pt]}, line width=1.1pt},
   body/.style={line width=0.6pt, fill=black!8}]
   % the cart, on the left
   \draw[body] (-0.7,-0.4) rectangle (0.7,0.4);
   \node at (0,0) {\small cart};
   \draw[force, accent] (0.7,0.1) -- (2.0,0.1)
      node[midway, above=0.3cm] {\small horse on cart};
   \draw[force] (-0.7,-0.25) -- (-1.2,-0.25)
      node[left] {\small ground on wheels};
   \draw[force] (0,0.4) -- (0,1.1) node[right] {\small normal force};
   \draw[force] (0,-0.4) -- (0,-1.1) node[right] {\small weight};
   % the horse, on the right
   \begin{scope}[xshift=5.0cm]
      \draw[body] (-0.7,-0.4) rectangle (0.7,0.4);
      \node at (0,0) {\small horse};
      \draw[force, accent] (-0.7,0.1) -- (-2.0,0.1)
         node[midway, above=0.3cm] {\small cart on horse};
      \draw[force] (0.7,-0.25) -- (2.2,-0.25)
         node[right] {\small ground on hooves};
      \draw[force] (0,0.4) -- (0,1.1) node[right] {\small normal force};
      \draw[force] (0,-0.4) -- (0,-1.1) node[right] {\small weight};
   \end{scope}
\end{tikzpicture}
\end{center}
The two forces of the third-law pair act on
different bodies, the horse's pull on the cart and the cart's pull on the
horse, so they never cancel in the second law of either. Whether each
body speeds up depends on all the forces on that body. On the cart, the
horse's forward pull is larger than the backward friction at its wheels,
so the cart accelerates forwards. On the horse, the ground pushes forwards
on the hooves, the third-law partner of the hooves' backward push on the
ground,
and this push is larger than the cart's backward pull, so the horse
accelerates forwards too. Taking horse and cart together, their pulls on
each other are internal and cancel (\cref{sec:ne-bodies}); the pair speeds
up because the forward push of the ground on the hooves exceeds the
backward friction on the wheels.

\solutionto{ex:ne-rope}The free-body diagram of the box has two forces,
the tension $F_{\mathrm{rope}}$ upwards and the weight $mg$ downwards,
the tension drawn longer because the box accelerates upwards.
\begin{center}
\begin{tikzpicture}[scale=1.25, line cap=round,
   force/.style={-{Stealth[length=6pt]}, line width=1.2pt},
   body/.style={line width=0.6pt, fill=black!8}]
   \draw[line width=0.8pt] (0,0.3) -- (0,1.9);
   \draw[body] (-0.4,-0.3) rectangle (0.4,0.3);
   \draw[force] (0,0) -- (0,1.45)
      node[midway, right=2pt] {\small $F_{\mathrm{rope}}$};
   \draw[force] (0,0) -- (0,-1.15) node[right] {\small $mg$};
   \fill (0,0) circle (1.4pt);
   \draw[accent, -{Stealth[length=5pt]}, dashed, line width=0.8pt]
      (1.3,-0.4) -- (1.3,0.5) node[right] {\small $a$};
\end{tikzpicture}
\end{center}
Taking upwards as positive, the second law gives
\begin{equation*}
   F_{\mathrm{rope}} - mg = ma
   \quad\Longrightarrow\quad
   F_{\mathrm{rope}} = m(g + a) = 2\times(9.80665 + 1.5)
   = \SI{22.6}{\newton} .
\end{equation*}
The rope must hold up the weight, \SI{19.6}{\newton}, and supply a further
$ma = \SI{3}{\newton}$ to accelerate the box.

\solutionto{ex:ne-moon}Mass measures inertia and is the same everywhere,
so the astronaut's mass on the Moon is \SI{70}{\kilogram}. The weight is
$mg$ with the Moon's value of $g$, $70\times1.62 = \SI{113}{\newton}$. On
the Earth it is $70\times9.80665 = \SI{686}{\newton}$, so on the Moon the
astronaut weighs $1.62/9.80665 = 0.165$ of their weight on the Earth, about
a sixth.

\solutionto{ex:ne-state}Measuring the height $x$ upwards from the ground,
the state at $t = 0$ is $x_0 = \SI{3}{\metre}$, $v_0 =
\SI{-2}{\metre\per\second}$, and by \cref{eq:mo-constant} with $a = -g$,
\begin{equation*}
   x(t) = 3 - 2t - \tfrac12 g t^2 .
\end{equation*}
The ball reaches the ground when $x = 0$. Multiplying by $-1$ to make the
leading coefficient positive, $\tfrac12 g t^2 + 2t - 3 = 0$, a quadratic
$At^2 + Bt + C = 0$ with $A = \tfrac12 g = 4.9033$, $B = 2$ and $C = -3$. Its roots are
$t = [-B \pm (B^2 - 4AC)^{1/2}]/2A$, and the positive one is
\begin{equation*}
   t = \frac{-2 + \sqrt{4 + 4\times4.9033\times3}}{2\times4.9033}
   = \frac{-2 + 7.927}{9.807} = \SI{0.604}{\second} ,
\end{equation*}
the red curve of \cref{fig:ne-state}(a).

\solutionto{ex:ne-euler}One step of length $\tau$ from $x = 0$, $v = v_0$
gives $x = v_0\tau$, while the true height is $v_0\tau - \tfrac12 g\tau^2$
by \cref{eq:mo-ball}. The step overshoots by $v_0\tau - (v_0\tau -
\tfrac12 g\tau^2) = \tfrac12 g\tau^2$. Two steps of $\tau/2$ go as
follows. The first gives $x_1 = v_0\tau/2$ and $v_1 = v_0 - g\tau/2$; the
second starts from this state, so
\begin{align*}
   x_2 &= x_1 + v_1\,\frac{\tau}{2}
      = \frac{v_0\tau}{2} + \Bigl(v_0 - \frac{g\tau}{2}\Bigr)\frac{\tau}{2}\\
   &= v_0\tau - \tfrac14 g\tau^2
      && \text{expanding; the two $v_0\tau/2$ add.}
\end{align*}
The overshoot is now $(v_0\tau - \tfrac14 g\tau^2) - (v_0\tau - \tfrac12
g\tau^2) = \tfrac14 g\tau^2$, half that of the single step. Halving the
step halves the error over a fixed time, though each step's own error
falls by four.

\solutionto{ex:ne-above}With $\gamma = \SI{5}{\per\second}$ the terminal
velocity is $v_\infty = 9.80665/5 = \SI{1.961}{\metre\per\second}$, and
the starting gap is $v_0 - v_\infty = 4 - 1.961 =
\SI{2.039}{\metre\per\second}$. At $t = \SI{0.2}{\second}$, $\gamma t = 1$,
so by \cref{eq:ne-drag-v} the gap has shrunk by the factor $\mathrm{e}^{-1}
= 0.368$:
\begin{equation*}
   v = 1.961 + 2.039\times0.368 = 1.961 + 0.750
   = \SI{2.711}{\metre\per\second} .
\end{equation*}
The bead slows towards $v_\infty$ from above.

\solutionto{ex:ne-approach}From rest, $v = v_\infty(1 -
\mathrm{e}^{-\gamma t})$, so \SI{99}{\percent} of $v_\infty$ is reached
when $\mathrm{e}^{-\gamma t} = 0.01$. Taking the natural logarithm of both
sides, $-\gamma t = \ln 0.01 = -4.605$, so $t = 4.605/5 =
\SI{0.921}{\second}$. Each further factor of ten closer to $v_\infty$
needs $\mathrm{e}^{-\gamma t}$ to shrink by another factor of 10. Since
$\mathrm{e}^{-\gamma(t + \Delta t)} = \mathrm{e}^{-\gamma t}\,
\mathrm{e}^{-\gamma\Delta t}$, this takes a further $\Delta t$ with
$\mathrm{e}^{-\gamma\Delta t} = 0.1$, that is $\Delta t = 2.303/\gamma =
\SI{0.461}{\second}$, as in \cref{sec:ne-drag}.

\solutionto{ex:ne-drag-check}Differentiating \cref{eq:ne-drag-x} term by
term, the derivative of $v_\infty t$ is $v_\infty$, and the derivative of
$1 - \mathrm{e}^{-\gamma t}$ is $\gamma\,\mathrm{e}^{-\gamma t}$ by the
chain rule, so
\begin{equation*}
   \frac{\dd x}{\dd t} = v_\infty - \frac{v_\infty}{\gamma}\,
   \gamma\,\mathrm{e}^{-\gamma t} = v_\infty\bigl(1 - \mathrm{e}^{-\gamma
   t}\bigr) ,
\end{equation*}
the velocity of a bead released from rest. Differentiating again,
$\dd^2 x/\dd t^2 = \gamma v_\infty\,\mathrm{e}^{-\gamma t}$. The right
side of \cref{eq:ne-drag} is
\begin{align*}
   g - \gamma\,\frac{\dd x}{\dd t}
   &= g - \gamma v_\infty + \gamma v_\infty\,\mathrm{e}^{-\gamma t}
      && \text{multiplying out the bracket}\\
   &= \gamma v_\infty\,\mathrm{e}^{-\gamma t}
      && \text{since $\gamma v_\infty = g$,}
\end{align*}
which equals $\dd^2 x/\dd t^2$, as required.

\solutionto{ex:ne-spring}By \cref{eq:ne-spring-omega}, $\omega =
\sqrt{k/m} = \sqrt{80/0.2} = \sqrt{400} = \SI{20}{\radian\per\second}$ and
the period is $2\pi/\omega = 2\pi/20 = \SI{0.314}{\second}$. Released from
rest at \SI{3}{\centi\metre}, the amplitude is $A = \SI{0.03}{\metre}$.
The largest speed, at the centre, is $A\omega = 0.03\times20 =
\SI{0.6}{\metre\per\second}$, and the largest acceleration, at the turning
points, is $A\omega^2 = 0.03\times400 = \SI{12}{\metre\per\second
\squared}$.

\solutionto{ex:ne-kick}Here $x_0 = 0$ and $v_0 = \SI{0.5}{\metre\per
\second}$, so
\begin{equation*}
   A = \sqrt{x_0^2 + (v_0/\omega)^2} = \sqrt{0 + 0.05^2}
   = \SI{0.05}{\metre} ,
\end{equation*}
with $\cos\phi = x_0/A = 0$ and $\sin\phi = -v_0/(\omega A) =
-0.5/(10\times0.05) = -1$. The angle with cosine 0 and sine $-1$ is $\phi =
-\pi/2$. The motion is therefore $x = 0.05\cos(10t - \pi/2)$. By
$\cos(-\theta) = \cos\theta$ and $\cos(\pi/2 - \theta) = \sin\theta$ from
\cref{sec:mo-trig}, $\cos(\omega t - \pi/2) = \cos(\pi/2 - \omega t) =
\sin\omega t$, so $x = 0.05\sin 10t$: the block starts at the centre and
moves outwards.

\solutionto{ex:ne-hanging}Let the spring be stretched by $s$ when the
block hangs at rest. The spring then pulls up with $ks$ and the weight
pulls down with $mg$, and since the block does not accelerate, $ks = mg$,
so $s = mg/k$. Now let the block be a distance $x$ below this resting
place, so that the spring is stretched by $s + x$. Taking downwards as
positive, the net force is
\begin{align*}
   mg - k(s + x) &= (mg - ks) - kx
      && \text{multiplying out the bracket}\\
   &= -kx
      && \text{since $ks = mg$,}
\end{align*}
and the second law gives $m\,\dd^2 x/\dd t^2 = -kx$, or $\dd^2 x/\dd t^2 =
-(k/m)\,x$. This is \cref{eq:ne-spring}, so the block oscillates about its
new resting place with the same $\omega = \sqrt{k/m}$ and period as on the
table. Gravity only moves the centre down by $mg/k$.

\solutionto{ex:ne-coupling}On the level, smooth track no external force
acts along the track, so the total momentum is the same before and after
the carts couple:
\begin{equation*}
   2\times3 + 1\times0 = (2 + 1)\,v
   \quad\Longrightarrow\quad
   v = \frac{6}{3} = \SI{2}{\metre\per\second} .
\end{equation*}

\solutionto{ex:ne-boat}No horizontal external force acts on the person
and boat together, and they start at rest, so their centre of mass stays
where it is. Let the person move $\Delta x_{\mathrm p}$ and the boat
$\Delta x_{\mathrm b}$ relative to the water, positive in the direction
the person walks. The centre of mass is at $X = (70\,x_{\mathrm p} +
140\,x_{\mathrm b})/210$, and when the person moves by $\Delta
x_{\mathrm p}$ and the boat by $\Delta x_{\mathrm b}$ it moves by $(70\,
\Delta x_{\mathrm p} + 140\,\Delta x_{\mathrm b})/210$. Keeping it fixed
therefore requires $70\,\Delta x_{\mathrm p} + 140\,\Delta x_{\mathrm b} =
0$. Walking \SI{3}{\metre}
along the boat means $\Delta x_{\mathrm p} = \Delta x_{\mathrm b} + 3$.
Substituting,
\begin{align*}
   70(\Delta x_{\mathrm b} + 3) + 140\,\Delta x_{\mathrm b} &= 0\\
   70\,\Delta x_{\mathrm b} + 210 + 140\,\Delta x_{\mathrm b} &= 0
      && \text{multiplying out}\\
   210\,\Delta x_{\mathrm b} &= -210
      && \text{collecting $\Delta x_{\mathrm b}$, subtracting 210}\\
   \Delta x_{\mathrm b} &= \SI{-1}{\metre}
      && \text{dividing by 210.}
\end{align*}
The boat moves \SI{1}{\metre} backwards, and the person \SI{2}{\metre}
forwards, relative to the water.

\solutionto{ex:ne-oxygen}By \cref{sec:ne-atoms}, a force in
\si{\electronvolt\per\angstrom} divided by a mass in \si{\amu} and
multiplied by \num{9.648533e-3} gives the acceleration in
\si{\angstrom\per\femto\second\squared}:
\begin{equation*}
   a = \frac{0.5}{15.999}\times9.648533\times10^{-3}
   = \SI{3.015e-4}{\angstrom\per\femto\second\squared} .
\end{equation*}
From rest, after \SI{5}{\femto\second} the velocity is $at = 5\times3.015
\times10^{-4} = \SI{1.508e-3}{\angstrom\per\femto\second}$. Since
$\SI{1}{\angstrom\per\femto\second} = \SI{e-10}{\metre}/\SI{e-15}{\second}
= \SI{e5}{\metre\per\second}$, this is \SI{151}{\metre\per\second}.

\solutionto{ex:ne-drift}Write $\vb{u} = \vb{P}/M$ for the drift and
$\vb{v}_i' = \vb{v}_i - \vb{u}$ for the new velocities. The new total
momentum is
\begin{align*}
   \sum_i m_i\vb{v}_i'
   &= \sum_i m_i\bigl(\vb{v}_i - \vb{u}\bigr)
      && \text{substituting $\vb{v}_i'$}\\
   &= \sum_i m_i\vb{v}_i - \sum_i m_i\vb{u}
      && \text{multiplying out each term}\\
   &= \sum_i m_i\vb{v}_i - \Bigl(\sum_i m_i\Bigr)\vb{u}
      && \text{$\vb{u}$ is common to every term}\\
   &= \vb{P} - M\,\frac{\vb{P}}{M} = \vb{0} .
\end{align*}
Subtracting the same constant $\vb{u}$ from every velocity is the change
to a frame of reference moving at $\vb{u}$ (\cref{sec:ne-frames}), in
which $\vb{a}' = \vb{a}$, so no acceleration changes. Every atom's
position is shifted by the same amount $\vb{u}t$, so the distances
between atoms, on which their forces depend, are unchanged as well.
